\section{Introduction}
    Digital images are conventionally stored and processed as three-dimensional arrays of shape $H \times W \times C$, where $H$ and $W$ denote the spatial dimensions and $C$ the number of color channels. Although matrix-based compression methods such as truncated Singular Value Decomposition (SVD) are well-established, they require either flattening the channel axis or processing each channel independently, thereby discarding inter-channel correlations that carry perceptually significant structure.\\

    Tensor decomposition methods naturally preserve this multi-way structure. A \emph{tensor} $\mathcal{X} \in \mathbb{R}^{I_1 \times I_2 \times \cdots \times I_N}$ generalizes matrices to $N$ modes, and \emph{low-rank tensor approximation} seeks a compact representation $\hat{\mathcal{X}}$ such that

    \begin{equation}
        \hat{\mathcal{X}} \;=\; \arg\min_{\hat{\mathcal{X}} \in \mathcal{M}_r}
        \mathcal{L}\!\left(\mathcal{X},\, \hat{\mathcal{X}}\right),
        \label{eq:general_problem}
    \end{equation}

    \noindent where $\mathcal{M}_r$ denotes a suitable low-rank manifold and $\mathcal{L}$ is a reconstruction loss. The choice of $\mathcal{L}$ is consequential: the squared Frobenius norm $\|\cdot\|_F^2$ yields analytically tractable subproblems and is the standard baseline, yet it is sensitive to outliers and heavy-tailed noise. Alternative losses---the $\ell_1$ norm and the Huber loss---offer robustness at the cost of algorithmic complexity, motivating a comparative study.\\

    \subsection*{Research Questions}
        This project investigates the following questions:

        \begin{enumerate}
            \item \textbf{Decomposition structure.}
                How do Tucker decomposition and CP decomposition differ in their approximation quality, compression ratio, and computational cost for image tensors?
            \item \textbf{Norm dependence.}
                How does the choice of reconstruction norm ($\|\cdot\|_F^2$, $\|\cdot\|_1$, Huber loss) affect the recovered factor matrices and the visual quality of the reconstructed image, particularly under structured noise?
            \item \textbf{Optimization behavior.}
                How do Alternating Least Squares (ALS) / Higher-Order Orthogonal Iteration (HOOI) compare with gradient-based methods (Adam, L-BFGS) in terms of convergence speed, stability, and sensitivity to initialization?
            \item \textbf{Rank--quality tradeoff.}
                Can the theoretical error bound derived from the multilinear singular value spectrum predict the empirically observed reconstruction error as a function of Tucker rank?
        \end{enumerate}

    \subsection*{Contributions}
        The contributions of this work are as follows.

        \begin{itemize}
            \item We implement Tucker decomposition from scratch (HOSVD initialization and HOOI refinement) and verify correctness against TensorLy \cite{kossaifi2019tensorly}.
            \item We derive and implement an Iteratively Reweighted Least Squares (IRLS) algorithm for $\ell_1$-Tucker and compare it with the Frobenius baseline and Huber-loss gradient descent on noisy image tensors.
            \item We conduct controlled numerical experiments on rank truncation, norm sensitivity, convergence, and initialization, connecting each empirical result to its theoretical prediction.
            \item We demonstrate a lightweight image compression pipeline that reports PSNR and compression ratio as functions of Tucker rank, providing an interpretable application of the mathematical framework.
        \end{itemize}

    \subsection*{Paper Organization}
        The remainder of this report is organized as follows. Section~\ref{sec:background} establishes the mathematical background, including tensor notation, Tucker and CP decomposition, and the norms used throughout. Section~\ref{sec:optimization} derives the optimization algorithms: ALS/HOOI for the Frobenius objective and IRLS plus gradient-based methods for robust objectives. Section~\ref{sec:norms} analyzes the statistical and geometric implications of each norm choice. Section~\ref{sec:experiments} presents numerical experiments that validate the theoretical predictions. Section~\ref{sec:demo} applies the framework to image compression. Section~\ref{sec:discussion} interprets the results and discusses limitations, and Section~\ref{sec:conclusion} concludes.\\

% % ============================================================
% %  sections/chapter_01_ko.tex
% %  1장: 서론
% % ============================================================

% \section{서론}
% \label{sec:introduction}

% 디지털 이미지는 일반적으로 $H \times W \times C$ 형태의 3차원 배열로
% 저장되고 처리된다. 여기서 $H$와 $W$는 공간 차원, $C$는 색상 채널 수를 나타낸다.
% 행렬 기반 압축 방법인 절단 특이값 분해(Truncated SVD)는 이미 잘 정립된
% 기법이지만, 채널 축을 평탄화(flatten)하거나 각 채널을 독립적으로 처리해야 하므로
% 지각적으로 중요한 구조를 담고 있는 채널 간 상관관계를 버리게 된다.

% 텐서 분해 방법은 이러한 다중 방향(multi-way) 구조를 자연스럽게 보존한다.
% 텐서 $\tensor{X} \in \mathbb{R}^{I_1 \times I_2 \times \cdots \times I_N}$은
% 행렬을 $N$개의 모드(mode)로 일반화한 것이며,
% \emph{저랭크 텐서 근사(low-rank tensor approximation)}는
% 다음과 같은 간결한 표현 $\hat{\tensor{X}}$를 구하는 문제이다:
% \begin{equation}
%   \hat{\tensor{X}}
%   \;=\; \arg\min_{\hat{\tensor{X}} \in \mathcal{M}_r}
%         \mathcal{L}\!\left(\tensor{X},\, \hat{\tensor{X}}\right),
%   \label{eq:general_problem}
% \end{equation}
% 여기서 $\mathcal{M}_r$은 적절한 저랭크 다양체(low-rank manifold)이고
% $\mathcal{L}$은 재구성 손실 함수이다.

% $\mathcal{L}$의 선택은 결과에 중대한 영향을 미친다.
% 프로베니우스 제곱 노름 $\|\cdot\|_F^2$은 해석적으로 다루기 쉬운
% 부분 문제를 생성하며 표준 기준선으로 널리 쓰인다.
% 그러나 이상치(outlier)와 중꼬리(heavy-tailed) 잡음에 민감하다는 단점이 있다.
% 반면 $\ell_1$ 노름과 Huber 손실은 강인성(robustness)을 제공하지만
% 알고리즘적 복잡성이 증가하므로, 이에 대한 비교 연구가 필요하다.

% % ------------------------------------------------------------
% \subsection*{연구 질문}
% % ------------------------------------------------------------

% 본 프로젝트는 다음 네 가지 질문을 탐구한다.

% \begin{enumerate}
%   \item \textbf{분해 구조.}
%         Tucker 분해와 CP 분해는 이미지 텐서에 대해
%         근사 품질, 압축률, 계산 비용 면에서 어떻게 다른가?

%   \item \textbf{노름 의존성.}
%         재구성 노름($\|\cdot\|_F^2$, $\|\cdot\|_1$, Huber 손실)의 선택이
%         복원된 인수 행렬과 이미지의 시각적 품질에, 특히 구조적 잡음 하에서,
%         어떤 영향을 미치는가?

%   \item \textbf{최적화 거동.}
%         교대 최소제곱법(ALS)/고차 직교 반복(HOOI)과
%         경사 기반 방법(Adam, L-BFGS)은 수렴 속도, 안정성,
%         초기화 민감도 측면에서 어떻게 비교되는가?

%   \item \textbf{랭크--품질 트레이드오프.}
%         다중선형 특이값 스펙트럼으로 유도한 이론적 오차 상한이
%         Tucker 랭크의 함수로서 실험적으로 관측된 재구성 오차를
%         얼마나 정확하게 예측하는가?
% \end{enumerate}

% % ------------------------------------------------------------
% \subsection*{기여}
% % ------------------------------------------------------------

% 본 보고서의 기여는 다음과 같다.

% \begin{itemize}
%   \item HOSVD 초기화와 HOOI 정제를 포함한 Tucker 분해를 처음부터 구현하고,
%         TensorLy~\cite{kossaifi2019tensorly}와 비교하여 정확성을 검증한다.

%   \item $\ell_1$-Tucker를 위한 반복 재가중 최소제곱(IRLS) 알고리즘을 유도 및
%         구현하고, 잡음이 포함된 이미지 텐서에서 프로베니우스 기준선 및
%         Huber 손실 경사 하강법과 비교한다.

%   \item 랭크 절단, 노름 민감도, 수렴, 초기화에 관한 통제된 수치 실험을 수행하고,
%         각 실험 결과를 해당 이론적 예측과 연결한다.

%   \item Tucker 랭크의 함수로서 PSNR과 압축률을 보고하는 경량 이미지 압축
%         파이프라인을 시연하여 수학적 프레임워크의 해석 가능한 응용을 제시한다.
% \end{itemize}

% % ------------------------------------------------------------
% \subsection*{보고서 구성}
% % ------------------------------------------------------------

% 이후 장의 구성은 다음과 같다.
% \ref{sec:background}장에서는 텐서 표기법, Tucker 및 CP 분해,
% 그리고 본 연구에서 사용되는 노름을 포함한 수학적 배경을 정립한다.
% \ref{sec:optimization}장에서는 최적화 알고리즘을 유도한다:
% 프로베니우스 목적 함수를 위한 ALS/HOOI와
% 강인 목적 함수를 위한 IRLS 및 경사 기반 방법을 다룬다.
% \ref{sec:norms}장에서는 각 노름 선택의 통계적·기하학적 함의를 분석한다.
% \ref{sec:experiments}장에서는 이론적 예측을 검증하는 수치 실험을 제시하고,
% \ref{sec:demo}장에서는 이미지 압축에 프레임워크를 적용한다.
% \ref{sec:discussion}장에서는 결과를 해석하고 한계를 논의하며,
% \ref{sec:conclusion}장에서 결론을 맺는다.